\documentclass[10pt,a4paper]{article}
\usepackage[utf8]{inputenc}
\usepackage[T1]{fontenc}
\usepackage[brazil]{babel}
\usepackage{amsmath,amssymb}
\usepackage{geometry}
\usepackage{multicol}
\usepackage{enumitem}
\usepackage{tikz}
\usepackage{qrcode}
\usepackage{fancyhdr}
\usepackage{array}
\usepackage{graphicx}
\usetikzlibrary{calc}
\geometry{top=1.15cm,bottom=1.25cm,left=1.25cm,right=1.25cm}
\setlength{\parindent}{0pt}
\setlength{\columnsep}{0.65cm}
\setlist[enumerate]{leftmargin=*,itemsep=1pt,topsep=2pt}
\newenvironment{alts}{\begin{enumerate}[label=\Alph*),leftmargin=5.5mm]}{\end{enumerate}}
\newcommand{\qtitle}[1]{\medskip\noindent\textbf{Questão #1.}\ }
\newcommand{\bubble}[1]{\tikz[baseline=-0.6ex]\draw (0,0) circle (1.7mm) node {\scriptsize #1};}
\pagestyle{fancy}
\fancyhf{}
\fancyhead[L]{\small IFSP -- Câmpus Hortolândia}
\fancyhead[R]{\small Física}
\fancyfoot[C]{\thepage}
\begin{document}
\thispagestyle{empty}

% Marcadores geométricos da folha OMR
\begin{tikzpicture}[remember picture,overlay]
\fill ($(current page.north west)+(8mm,-8mm)$) rectangle ++(5mm,-5mm);
\fill ($(current page.north east)+(-13mm,-8mm)$) rectangle ++(5mm,-5mm);
\fill ($(current page.south west)+(8mm,13mm)$) rectangle ++(5mm,-5mm);
\fill ($(current page.south east)+(-13mm,13mm)$) rectangle ++(5mm,-5mm);
\end{tikzpicture}

\begin{center}
{\Large\bfseries <<TITLE>>}\\[2mm]
{\small <<SUBTITLE>>}
\end{center}

\vspace{2mm}
\noindent\textbf{Nome completo:}\ \rule{12.5cm}{0.4pt}\\[3mm]
\textbf{Turma:}\ \rule{4cm}{0.4pt}\hfill
\textbf{Data:}\ \rule{3.5cm}{0.4pt}

\vspace{4mm}
\begin{center}
\begin{minipage}{0.84\textwidth}
\centering
{\large\bfseries Folha de respostas}\\[2mm]
Preencha completamente uma única alternativa por questão com caneta azul ou preta.\\[4mm]
\begin{tabular}{c *{5}{>{\centering\arraybackslash}p{1.5cm}}}
\textbf{Q} & \textbf{A}&\textbf{B}&\textbf{C}&\textbf{D}&\textbf{E}\\[1mm]
<<OMR_ROWS>>
\end{tabular}

\vspace{5mm}
\begin{tabular}{p{0.58\linewidth}p{0.32\linewidth}}
\textbf{Código de controle:} <<PUBLIC_CODE>> &
\centering\qrcode[height=2.0cm]{<<QR_PAYLOAD>>}
\end{tabular}
\end{minipage}
\end{center}

\vfill
{\footnotesize O QR code identifica apenas a versão da prova; ele não contém o gabarito.}
\newpage

\begin{center}
{\large\bfseries <<TITLE>>}
\end{center}
<<FORMULA_SHEET>>

\begin{multicols}{2}
<<QUESTIONS>>
\end{multicols}
\end{document}
